\section{Methodology: The Aero Architecture}
Aero (Gen 6) addresses the distribution sensitivity of previous generations by replacing static binning with a quantile-adaptive mapping function. While linear range sorters suffer from bin-collisions on non-uniform data, Aero ensures balanced occupancy by approximating the dataset's distribution via a 1024-entry lookup table $\mathcal{L}$.

This is implemented through a 128-bit reciprocal fixed-point multiplier $\mathcal{R}$ calculated from a small sorted sample ($m=256$). The mapping $\mathcal{B}(x) = \mathcal{L} [ \lfloor (\mathcal{K}(x) - \min) \cdot \mathcal{R} \cdot 2^{-64} \rfloor ]$ resolves the bin index in a few CPU cycles. By limiting $\mathcal{L}$ to 2KB, we ensure the structure remains pinned in the L1 Data Cache throughout the scatter phase. Algorithm \ref{alg:aero_mapping} details this classification logic.

\begin{algorithm}
\caption{Aero Division-Free Quantile Mapping}
\label{alg:aero_mapping}
\begin{algorithmic}[1]
\REQUIRE Input key $k$, Min value $min$, Multiplier $R$, Table $L$
\ENSURE Bin index $b$
\STATE $diff \leftarrow k - min$
\STATE $idx \leftarrow (diff \cdot R) \gg 64$ 
\IF{isUniform}
    \STATE $b \leftarrow \min(idx, numBins - 1)$
\ELSE
    \STATE $tIdx \leftarrow \min(idx, 1023)$
    \STATE $b \leftarrow L[tIdx]$ 
\ENDIF
\RETURN $b$
\end{algorithmic}
\end{algorithm}

To maximize throughput, ARS Aero currently utilizes a scratch-buffer scatter. While in-place permutation is theoretically possible, the "Move Penalty" incurred by cycle-following and atomic swap-holes often outweighs the memory savings. By utilizing an $O(N)$ auxiliary buffer, we ensure writes are contiguous and unidirectional, allowing the hardware prefetchers to operate at peak efficiency.

\subsection{Stability Guarantees}
A critical requirement for enterprise systems is \textit{Stability}—the preservation of the relative order of elements with equal keys. ARS achieves this through a two-stage inheritance model.

\subsubsection{Formalization of Stability Inheritance}
Let $S = (x_1, x_2, \dots, x_n)$ be the input sequence and $\mathcal{B}(x)$ be the mapping function. For any two elements $x_i, x_j \in S$ where $x_i = x_j$ and $i < j$, the partitioning phase ensures:
\begin{equation}
\text{addr}(x_i) < \text{addr}(x_j) \text{ in } Bin_{\mathcal{B}(x_i)}
\end{equation}
This is guaranteed because the scatter pass traverses $S$ linearly, appending elements to their respective bins in their original arrival order. Consequently, global stability is satisfied if the intra-bin refiner $f_{refine}$ is also stable:
\begin{equation}
\text{ARS}_{stable}(S) = \bigoplus_{k=0}^{B-1} f_{stable}(Bin_k)
\end{equation}
where $\bigoplus$ denotes the ordered concatenation of partitioned runs.

\subsubsection{Architectural Divergence: Unstable vs. Stable Refinement}
The ARS framework allows for the dynamic substitution of the refinement engine based on the required invariants.
\begin{itemize}
    \item \textbf{Unstable Refinement (PDQsort)}: This approach utilizes Pattern-Defeating Quicksort, which optimizes for instruction-level parallelism by minimizing branches. However, its recursive partitioning can lead to non-sequential memory access within the L2 cache boundaries.
    \item \textbf{Stable Refinement (Timsort)}: This approach utilizes a merge-driven logic. While traditionally considered slower due to higher move counts, Timsort excels in the ARS environment because the bins are highly partitioned ($n_j \approx 1000$). At this scale, Timsort's sequential merge patterns maximize cache-line reuse and hardware prefetching, often yielding the superior 8.57 ms latency observed in our evaluations.
\end{itemize}


